\section*{부록: 자연 27장면의 궤적 판정 사유}
\label{sec:appendix-trajectory}
\begin{revision}
표~\ref{tab:trajectory-details}는 \ref{sec:trajectory-results}절의 장면별 궤적 판정과 사건 수준의 사유 건수다. 한 장면에 여러 사유가 겹칠 수 있으며 장면 판정은 \ref{sec:aux-design}절의 우선순위로 집계했다.

\begin{table*}[h]
\centering
\color{revisionblue}
\bitablecaption{27장면의 궤적 판정과 사건 사유 건수. 일치: 행동 일치, 불일치: 종방향 불일치, 해석: 구문 해석 미상, 미지원: 지원 행동 부재, 관측: 관측 행동 불명확, 시점: 단계 시점ㆍ해제 근거 부족.}{Trajectory verdicts and event-level reason counts for the 27 natural scenes.}
\label{tab:trajectory-details}
\scriptsize
\setlength{\tabcolsep}{3.5pt}
\begin{tabular}[t]{llrrrrrr}
\toprule
장면 & 판정 & 일치 & 불일치 & 해석 & 미지원 & 관측 & 시점 \\
\midrule
G01 & 미상 & 0 & 0 & 0 & 4 & 2 & 0 \\
G02 & 미상 & 0 & 0 & 2 & 1 & 0 & 0 \\
G03 & 불일치 & 1 & 1 & 2 & 0 & 0 & 1 \\
G04 & 미상 & 1 & 0 & 1 & 1 & 0 & 1 \\
G05 & 불일치 & 0 & 1 & 1 & 4 & 0 & 0 \\
G06 & 미상 & 2 & 0 & 0 & 0 & 0 & 2 \\
G07 & 미상 & 3 & 0 & 1 & 0 & 0 & 0 \\
G08 & 미상 & 4 & 0 & 0 & 3 & 0 & 2 \\
G09 & 미상 & 3 & 0 & 0 & 1 & 0 & 2 \\
G10 & 불일치 & 0 & 1 & 0 & 1 & 0 & 0 \\
G11 & 미상 & 3 & 0 & 1 & 1 & 0 & 0 \\
G12 & 미상 & 3 & 0 & 1 & 0 & 0 & 2 \\
G13 & 불일치 & 1 & 1 & 1 & 0 & 0 & 1 \\
G14 & 미상 & 1 & 0 & 0 & 0 & 0 & 2 \\
\bottomrule
\end{tabular}
\hspace{1.5em}
\begin{tabular}[t]{llrrrrrr}
\toprule
장면 & 판정 & 일치 & 불일치 & 해석 & 미지원 & 관측 & 시점 \\
\midrule
G15 & 미상 & 0 & 0 & 0 & 4 & 0 & 1 \\
G16 & 일치 & 3 & 0 & 0 & 0 & 0 & 0 \\
G17 & 미상 & 0 & 0 & 0 & 2 & 0 & 1 \\
G18 & 일치 & 3 & 0 & 0 & 0 & 0 & 0 \\
G19 & 미상 & 0 & 0 & 0 & 2 & 0 & 1 \\
G20 & 미상 & 0 & 0 & 4 & 3 & 0 & 0 \\
G21 & 미상 & 0 & 0 & 0 & 2 & 0 & 0 \\
G22 & 미상 & 1 & 0 & 0 & 3 & 0 & 0 \\
G23 & 불일치 & 4 & 1 & 0 & 0 & 0 & 0 \\
G24 & 불일치 & 2 & 3 & 0 & 0 & 0 & 0 \\
G25 & 불일치 & 0 & 1 & 1 & 1 & 0 & 1 \\
G26 & 미상 & 0 & 0 & 0 & 6 & 0 & 0 \\
G27 & 미상 & 2 & 0 & 0 & 2 & 0 & 0 \\
\\
\bottomrule
\end{tabular}
\end{table*}
\end{revision}
